% Vietnamese translation for OI Public Library
% Source-File: IOI中国国家候选队论文/2008/Day2/1.程芃祺《计算几何中的二分思想》/计算几何中的二分思想.ppt
\documentclass[11pt]{article}
\usepackage{oipl-vn}

\newcommand{\Oh}{\mathcal{O}}

\title{Bản trình chiếu: Tư tưởng chia đôi trong hình học tính toán}
\author{Cheng Pengqi}
\date{2008}

\begin{document}
\maketitle

\origtitle{Binary-search thinking in computational geometry: slides}
\sourcepath{IOI China 2008 / Day 2 / Cheng Pengqi / slides.ppt}

\section{Mạch trình bày}

PPT gốc trích được các mốc chính: \textit{Simplified GSM Network},
\textit{Collecting Luggage}, \textit{Heliport} và \textit{Flight Safety}.
Bản ghi chú này đi theo đúng mạch ấy. Slide không nhằm giới thiệu một cấu
trúc hình học mới, mà nhấn mạnh một cách nghĩ: khi bài toán hình học trực
tiếp dẫn tới dựng hình, liệt kê trường hợp hoặc tính cực trị quá rườm rà, hãy
thử chuyển nó thành một bài toán phán định có tính đơn điệu, hoặc chia nhỏ
đối tượng liên tục cho đến khi phần còn lại đủ đơn giản.

Trong bốn ví dụ, ``chia đôi'' xuất hiện dưới nhiều dạng. Có lúc đó là chia
đôi một đoạn thẳng và cộng kết quả hai nửa; có lúc là chia đôi thời gian hoặc
bán kính; có lúc là chia đôi đáp án để biến bài toán tối ưu hóa thành kiểm
tra phủ. Điểm chung là tránh phải dựng toàn bộ cấu trúc hình học phức tạp.

\section{\textit{Simplified GSM Network}: chia đoạn thay cho dựng Voronoi}

Bài \textit{Simplified GSM Network} có các trạm phát sóng, các thành phố và
các tuyến đường nối thành phố. Mỗi vị trí trên mặt phẳng dùng trạm gần nhất.
Với một tuyến đường là đoạn thẳng, chi phí cần biết chính là số lần đoạn ấy
đi qua biên giữa hai miền trạm khác nhau. Nếu dùng mô hình kinh điển, đây là
bài toán dựng biểu đồ Voronoi rồi đếm số cạnh Voronoi bị đoạn thẳng cắt.

Slide ghi trực tiếp các ký hiệu \(w[l]=0\), \(w[l]=1\), \(l=l_1+l_2\) và
\(w[l]=w[l_1]+w[l_2]\). Ý nghĩa là: xét đoạn \(l\). Nếu hai đầu mút của
\(l\) cùng thuộc một trạm gần nhất, đoạn này không gây chuyển trạm, nên
\(w[l]=0\). Nếu hai đầu mút thuộc hai trạm khác nhau nhưng đoạn đã đủ ngắn,
ta coi đoạn ấy cắt đúng một biên Voronoi, nên đóng góp \(1\). Trường hợp còn
lại, chia \(l\) tại trung điểm thành \(l_1,l_2\) rồi cộng hai trọng số.

Lợi ích không nằm ở độ phức tạp lý thuyết của biểu đồ Voronoi, mà ở tính thực
dụng. Dựng Voronoi trong phòng thi dễ kéo theo nhiều ca suy biến và lỗi số
thực. Còn với chia đoạn, phần cần cài đặt chỉ là tìm trạm gần nhất của một
điểm và một thủ tục đệ quy có ngưỡng \(\varepsilon\). Sau khi có trọng số của
từng tuyến đường, phần còn lại là đường đi ngắn nhất thông thường.

\section{\textit{Collecting Luggage}: biến động thành tĩnh}

Trong \textit{Collecting Luggage}, vali di chuyển dọc theo biên của một đa
giác đơn, còn người xuất phát từ một điểm trong mặt phẳng với vận tốc lớn hơn.
Cần tìm thời điểm sớm nhất để người lấy được vali. Nếu cố mô phỏng trực tiếp
đường đi tới một điểm đang chuyển động, ta phải xử lý nhiều tình huống: đích
nằm trên cạnh nào, từ điểm hiện tại nhìn thấy đoạn nào, và đường ngắn nhất
qua các đỉnh ra sao.

Cách nhìn của slide là chia đôi thời gian \(t\). Với một \(t\) đã cho, vị trí
của vali là một điểm cố định trên biên đa giác. Khi đó bài toán trở thành:
khoảng cách đường đi ngắn nhất từ người tới điểm cố định này có không vượt
quá \(V_p t\) hay không? Ta có thể thêm điểm vali vào biên, tách cạnh tương
ứng, rồi tính đường ngắn nhất trong đồ thị nhìn thấy được hoặc dùng các giá
trị đã tiền xử lý tới các đỉnh.

Điều làm phép chia đôi hợp lệ là tính đơn điệu. Nếu người bắt được vali tại
thời điểm \(t\), thì với thời gian lớn hơn người cũng có thể đi tới không
muộn hơn vị trí tương ứng; ngược lại, nếu \(t\) chưa đủ thì mọi thời điểm nhỏ
hơn cũng chưa đủ. Bài viết chứng minh điều này bằng bất đẳng thức tam giác và
điều kiện \(V_l<V_p\). Vì thế một bài toán động được chuyển thành chuỗi bài
toán tĩnh, mỗi bài chỉ kiểm tra một thời điểm.

\section{\textit{Heliport}: chia bán kính để giảm số ca đại số}

\textit{Heliport} yêu cầu đặt một hình tròn lớn nhất bên trong một đa giác
trực giao. Nếu giải trực tiếp, tâm và bán kính của đường tròn tối ưu thường
được xác định bởi ba điều kiện tiếp xúc. Với cạnh ngang, cạnh dọc và đỉnh
đa giác, các trường hợp có thể được ký hiệu là \(P,H,V\). Ba điều kiện tạo ra
nhiều kiểu như \(PPP, PPH, PPV, PHV, HHV, HVV, PHH, PVV\). Dù số ca hữu hạn,
việc lập và cài các phương trình rất dễ sai.

PPT ghi các mốc \(d_1=d_2=d_3\), \(d_1=d_2=d_3=r\) và \(d_1=d_2=r\). Đây là
đúng chỗ chia đôi phát huy tác dụng. Nếu bán kính \(r\) đã biết, ta không còn
cần giải đồng thời ba điều kiện chưa biết; chỉ cần kiểm tra có tồn tại tâm
nằm trong đa giác và cách biên ít nhất \(r\) hay không. Các trường hợp đại số
giảm từ bộ ba xuống những quan hệ đôi như \(PP,PH,PV,HV\), hoặc có thể hiểu
tương đương là kiểm tra giao của các miền hợp lệ sau khi co đa giác theo bán
kính \(r\).

Tính đơn điệu rất rõ: nếu một bán kính \(r\) đặt được, mọi bán kính nhỏ hơn
cũng đặt được; nếu \(r\) không đặt được, mọi bán kính lớn hơn cũng không đặt
được. Vì thế ta chia đôi \(r\), dùng thủ tục phán định để cập nhật cận, và
tránh được phần lớn tính toán công thức dễ lỗi.

\section{\textit{Flight Safety}: tìm bằng cách chứng minh phủ}

Trong \textit{Flight Safety}, đường bay là một đường gấp khúc, các lục địa là
những đa giác không giao nhau. Cần tìm giá trị lớn nhất, trên toàn bộ đường
bay, của khoảng cách tới lục địa gần nhất. Tìm cực trị trực tiếp khó vì điểm
xa nhất có thể nằm trong lòng một đoạn bay, chịu ảnh hưởng đồng thời bởi
nhiều đa giác.

Thay vì đi tìm điểm cực trị, ta thử một đáp án \(d\). Mỗi lục địa được nở ra
thêm khoảng cách \(d\): cạnh sinh ra dải chữ nhật, đỉnh sinh ra cung tròn.
Nếu mọi điểm trên đường bay đều nằm trong hợp các miền đã nở, thì khoảng cách
lớn nhất không vượt quá \(d\); nếu còn đoạn nào ở ngoài, \(d\) quá nhỏ.

Để kiểm tra, cắt từng đoạn bay bởi các giao điểm với cạnh và cung tròn của
các miền nở. Trên mỗi đoạn con, trạng thái trong--ngoài không đổi, nên chỉ
cần kiểm tra một điểm đại diện như trung điểm. Khoảng cách \(d\) càng lớn thì
miền phủ càng rộng, vì vậy điều kiện có tính đơn điệu và có thể chia đôi.

\section{Tổng kết}

Các slide cho thấy chia đôi trong hình học không chỉ là tìm kiếm nhị phân
trên số thực. Nó là một cách thay đổi hình dạng bài toán:
\begin{itemize}
  \item chia nhỏ đoạn để tránh dựng cấu trúc toàn cục;
  \item chia thời gian để biến chuyển động thành kiểm tra tĩnh;
  \item chia bán kính hoặc khoảng cách để thay tối ưu hóa bằng phán định;
  \item dùng tính đơn điệu để điều khiển sai số số thực một cách rõ ràng.
\end{itemize}

Thông điệp thực dụng nhất là: trước khi lao vào các công thức hình học dài,
hãy hỏi liệu một giá trị ứng viên có kiểm tra được dễ hơn không. Nếu câu trả
lời là có và điều kiện đơn điệu, chia đôi thường giúp lời giải ngắn hơn, ít
ca suy biến hơn và dễ tin cậy hơn.

\end{document}
